
\chapter{JSON (2 of 2)}
\label{ch:JSON2}

% TODO: A big-ol' memory diagram of the entire "courses" nested structure.

In this chapter we'll cover some common techniques for dealing with JSON data.
Most often, what you end up doing is writing Python code that combs
systematically through the data, surgically extracting the parts you need for
analysis. These parts can then be assembled in NumPy arrays, Pandas
\texttt{DataFrame}s, or whatever, and probed as you're accustomed to.

\index{pretty printing@``pretty printing''}

\label{coursesJsonPretty}

Let's go back to the college courses example from p.\pageref{coursesJson},
pretty printed here:

\index{courses (JSON file)}
\index{json file@\texttt{.json} file}
\index{DATA 219}
\index{Lady Gaga}
\index{Rock, The}

\begin{Verbatim}[fontsize=\small,samepage=true,frame=lines,framesep=3mm]
[{'creds': 3,
  'instructor': 'Davies',
  'meetings': [{'location': {'bldg': 'Trinkle', 'room': 'B6'},
                'time': {'day': 'TR', 'time': '12:30pm'}}],
  'num': 219,
  'prefix': 'DATA',
  'roster': [{'email': 'lgaga@umw.edu', 'first': 'Lady', 'last': 'Gaga'},
             {'email': 'trock@umw.edu', 'first': 'The', 'last': 'Rock'},
             ...,
             ...]
 },
 ...
]
\end{Verbatim}

\index{[]@\texttt{[]} (boxies)}
\index{boxies (square brackets)}
\index{\{\}@\texttt{\{\}} (curlies)}
\index{curlies (curly braces)}
\index{list}
\index{dictionary}

The first brute fact to grasp is that whatever its internal complexities, this
whole thing is \textit{a list of dictionaries}. Not a list of lists, or a
dictionary of dictionaries, or a dictionary of lists, but: a list of
dictionaries. How do we know this? Because the first character is a boxie
(``\texttt{[}'') and the second one is a curly (``\texttt{\{}'').

\section{Printing a class roster}

As a first example of working with this data set, let's say we wanted to simply
print the full names of all the students in the 17th course in this list.
First, let's load the \texttt{.json} file into Python:

\begin{Verbatim}[fontsize=\small,samepage=true,frame=single,framesep=3mm]
import json

with open("courses.json") as f:
    c = json.load(f)
\end{Verbatim}

\index{loop}

Here we've called the outermost structure ``\texttt{c}'', which stands for
``courses'' but is shortened since we'll type it a lot. (Remember that
\texttt{c} is a list -- a very long and complicated list, but still a list.)

And then we do a deep think. Printing all the names clearly requires a loop.
But what exactly do we want to loop through? The answer, after scrutinizing the
above structure, is: we need to find the value of the \texttt{"roster"} key of
the 17th dictionary in the \texttt{c} list. That in turn will be a list, and it
will contain dictionaries, one for each student. (Pause before continuing, look
at the JSON contents from p.~\pageref{coursesJsonPretty}, and see if you agree
with every word of this paragraph.)

So, the loop we want is something like this:

\begin{Verbatim}[fontsize=\small,samepage=true,frame=single,framesep=3mm]
for student in c[16]['roster']: 
    print("{} {} is in {} {}.".format(student['first'], 
        student['last'], c[16]['prefix'], c[16]['num'])) 
\end{Verbatim}
\vspace{-.2in}

\begin{Verbatim}[fontsize=\small,samepage=true,frame=leftline,framesep=5mm,framerule=1mm]
Lady Gaga is in DATA 219.
The Rock is in DATA 219.
Obi-wan Kenobi is in DATA 219.
Amy Klobuchar is in DATA 219.
...
\end{Verbatim}

Success. (Don't worry if you don't get this loop right the first time; it often
takes some tweaking and refinement before it works.) Notice how this code
operates:

\begin{compactenum}

\item We dig fairly deep into \texttt{c} (which is a list), finding its 17th
element (at index 16 of course), which is a dictionary. Then we get the value
of that dictionary's \texttt{"roster"} key. As you can see above, that key's
value is a list.

\item Every time through the loop, the ``\texttt{student}'' variable is set
equal to the next element of this list. And that element, as you can see above,
is a dictionary. The dictionary has two items of interest: \texttt{"first"} and
\texttt{"last"}. We then print the values of these two keys, plus the course
prefix and number. (Note that the course prefix and number come directly from
the \texttt{c[16]} dictionary itself, whereas the first and last names come
from keys of the elements of the inner \texttt{"roster"} list that we're
iterating through.)

\end{compactenum}

\section{Printing \textit{all} the class rosters}

That printed the roster for one class. What about printing the roster for
\textit{all} the classes?

If you think about it, the structure is much the same -- except that instead of
merely iterating through class \#17's list of dictionaries, we want to do it
for \textit{every} class.

\index{nested!loops}

You'll probably realize that this necessitates \textit{another}. All the stuff
we did for class \#17, we now want to do for all the classes, one after
another. So in essence, we wrap the previous code in a nested, outer loop which
will perform it once for each course. Check out the following:

\begin{Verbatim}[fontsize=\small,samepage=true,frame=single,framesep=3mm]
for course in c:
    roster = course['roster']
    for student in roster:
        print("{} {} is in {} {}.".format(student['first'], 
            student['last'], course['prefix'], course['num'])) 
\end{Verbatim}

Observe the nested structure (a \texttt{for} within a \texttt{for}), and the
use of a temporary helper variable (\texttt{"roster"}) to make the code
clearer. Also notice how the \texttt{print()} statement includes information
from both the \texttt{student} and \texttt{course} variables, as appropriate.

This is just the right thing. However, when we run it, we get an error after a
while:

\begin{Verbatim}[fontsize=\small,samepage=true,frame=leftline,framesep=5mm,framerule=1mm]
Lady Gaga is in DATA 219.
The Rock is in DATA 219.
Obi-wan Kenobi is in DATA 219.
...
Elizabeth Jennings is in PSCI 401.
Stan Beeman is in PSCI 401.
Phillip Jennings is in PSCI 401.
...
Tyrion Lannister is in ENGL 101.
Jon Snow is in ENGL 101.
Daenerys Targaryen is in ENGL 101.
...
Han Solo is in CHEM 211.
Poe Dameron is in CHEM 211.
---------------------------------------------------------------------------
KeyError                              Traceback (most recent call last)
     77     for student in roster:
     78         print("{} {} is in {} {}.".format(student['first'], 
---> 79             student['last'], course['prefix'], course['num'])) 
     80 

KeyError: 'last'
\end{Verbatim}

Whoaaaa, what's \textit{dat}?

\subsection{Encountering erroneous/missing data}

This kind of error will happen all the time when you're working with JSON. The
first step is not to panic. The second step is to calmly observe that
apparently, everything was working fine up to a point...and then it suddenly
broke. The third step is to look at that error message and see if you can make
sense out of it. (Often, you can.) And the fourth step is to introduce some
tactical debugging messages into your code, so that you can narrow down the
source of the error.

I assume you've done those first two steps. On to the third: the error message.
It says, quote:

\begin{Verbatim}[fontsize=\normalsize,samepage=true,frame=leftline,framesep=5mm,framerule=1mm]
KeyError: 'last'
\end{Verbatim}

and it evidently happened on line 79, which was in the middle of the
\texttt{print()} statement. Put 2 and 2 together and you'll realize that when
we attempted to extract the value of the \texttt{"last"} key for one of the
dictionaries, \textit{it wasn't there.} That's what ``\texttt{KeyError}'' turns
out to mean: you asked for a particular key but there wasn't one in the
dictionary.

The fourth step is crucial. It marks the difference between fledgling, flailing
novices and cool, critical-thinking experts. We need to insert (temporarily)
one or more \textit{print()} statements into our loop, so that \textit{right
before the program crashes,} we get some crucial insight into what was about
to go wrong.

This takes some practice and some intuition, but in this case let me simply
observe what it makes sense to do. If we were about to access a non-existent
key from a dictionary, it might illuminate things if we printed out the
dictionary immediately before trying that stunt. That way, we'll see the state
of affairs just before the error.

Here's the key line, inserted into the inner loop:

\begin{Verbatim}[fontsize=\small,samepage=true,frame=single,framesep=3mm]
for course in c:
    roster = course['roster']
    for student in roster:
        print("The dict is: {}".format(student))
        print("{} {} is in {} {}.".format(student['first'], 
            student['last'], course['prefix'], course['num'])) 
\end{Verbatim}

Now look at the revised output, below:

\begin{Verbatim}[fontsize=\footnotesize,samepage=true,frame=leftline,framesep=5mm,framerule=1mm]
The dict is: {'email': 'han.solo@corellia.gov', 'first': 'Han', 'last': 'Solo'}
Han Solo is in CHEM 211.
The dict is: {'email': 'dameron@resistance.org', 'first': 'Poe', 'last': 'Dameron'}
Poe Dameron is in CHEM 211.
The dict is: {'email': 'root@dagobah.net', 'first': 'Yoda'}
---------------------------------------------------------------------------
KeyError                                  Traceback (most recent call last)
     78         print("The dict is: {}".format(student))
     79         print("{} {} is in {} {}.".format(student['first'], 
---> 80             student['last'], course['prefix'], course['num'])) 
     81 

KeyError: 'last'
\end{Verbatim}

\index{Yoda}
\index{CHEM 211}
\index{DATA 219}

Aha! Our problem is \texttt{Yoda}, a student in CHEM 211 who doesn't have a
last name.

It takes practice and experience to get a feel for where to insert meaningful
debugging statements. The key question to ask is: ``what do I wish I had known
right before the error occurred?'' Add lines of code to shed light on that
question.


\subsection{Dealing with erroneous/missing data}

Once you've discovered the absence of something you were looking for, you need
to ask yourself what the implications are of it being missing. Sometimes it's a
bona fide data error that means you need to investigate the source of the
information, discard some data points, or something else. Other times (as in
the Yoda case) it's innocuous -- not everyone has a last name, after all -- but
you still have to \textit{deal} with the problem somehow so your program can
continue.

In the case of a missing key which you can sensibly ignore, you can code around
it as follows:

\begin{Verbatim}[fontsize=\small,samepage=true,frame=single,framesep=3mm]
for course in c:
    roster = course['roster']
    for student in roster:
        if 'last' in student:
            print("{} {} is in {} {}.".format(student['first'], 
                student['last'], course['prefix'], course['num'])) 
        else:
            print("{} is in {} {}.".format(student['first'], 
                course['prefix'], course['num'])) 
\end{Verbatim}

\index{in@\texttt{in}}

The important line is the fourth one: ``\texttt{if \textquotesingle
last\textquotesingle~in student:}''. The ``\texttt{in}'' keyword checks to see
whether a dictionary has a certain key or not. In this case, if there is no
\texttt{"last"} key, we want to simply print the first name and go on, not
crash. The output is now satisfactory:

\begin{Verbatim}[fontsize=\small,samepage=true,frame=leftline,framesep=5mm,framerule=1mm]
...
Han Solo is in CHEM 211.
Poe Dameron is in CHEM 211.
Yoda is in CHEM 211.
Emperor Palpatine is in CHEM 211.
Chewie is in CHEM 211.
Rey is in CHEM 211.
...
\end{Verbatim}

The appropriate way to deal with an error like this varies depending on the
situation. This is why you need to back away from the keyboard, take a walk,
and ask yourself, ``what does it \textit{mean} when this value is absent or
corrupt? Does that imply that I need to throw out that data point? Or that
something deep and sinister is going on? Or is there a way to safely ignore
it?''
